%!TeX root=../gardentop.tex
\chapter{A Young Rajah} 
	
\begin{figure}[t!]
\centering
\includegraphics[width=\linewidth]{april}
\end{figure}

	\lettrine[]{T}{he} moor was hidden in mist when the morning came and the rain had not stopped pouring down. There could be no going out of doors. Martha was so busy that Mary had no opportunity of talking to her, but in the afternoon she asked her to come and sit with her in the nursery. She came bringing the stocking she was always knitting when she was doing nothing else.

\zz
<What's the matter with thee?> she asked as soon as they sat down. <Tha' looks as if tha'd somethin' to say.>

<I have. I have found out what the crying was,> said Mary.

Martha let her knitting drop on her knee and gazed at her with startled eyes.

<Tha' hasn't!> she exclaimed. <Never!>

<I heard it in the night,> Mary went on. <And I got up and went to see where it came from. It was Colin. I found him.>

Martha's face became red with fright.

<Eh! Miss Mary!> she said half crying. <Tha' shouldn't have done it—tha' shouldn't! Tha'll get me in trouble. I never told thee nothin' about him—but tha'll get me in trouble. I shall lose my place and what'll mother do!>

<You won't lose your place,> said Mary. <He was glad I came. We talked and talked and he said he was glad I came.>

<Was he?> cried Martha. <Art tha' sure? Tha' doesn't know what he's like when anything vexes him. He's a big lad to cry like a baby, but when he's in a passion he'll fair scream just to frighten us. He knows us daren't call our souls our own.>

<He wasn't vexed,> said Mary. <I asked him if I should go away and he made me stay. He asked me questions and I sat on a big footstool and talked to him about India and about the robin and gardens. He wouldn't let me go. He let me see his mother's picture. Before I left him I sang him to sleep.>

Martha fairly gasped with amazement.

<I can scarcely believe thee!> she protested. <It's as if tha'd walked straight into a lion's den. If he'd been like he is most times he'd have throwed himself into one of his tantrums and roused th' house. He won't let strangers look at him.>

<He let me look at him. I looked at him all the time and he looked at me. We stared!> said Mary.

<I don't know what to do!> cried agitated Martha. <If Mrs~Medlock finds out, she'll think I broke orders and told thee and I shall be packed back to mother.>

<He is not going to tell Mrs~Medlock anything about it yet. It's to be a sort of secret just at first,> said Mary firmly. <And he says everybody is obliged to do as he pleases.>

<Aye, that's true enough—th' bad lad!> sighed Martha, wiping her forehead with her apron.

<He says Mrs~Medlock must. And he wants me to come and talk to him every day. And you are to tell me when he wants me.>

<Me!> said Martha; < I shall lose my place—I shall for sure!>

<You can't if you are doing what he wants you to do and everybody is ordered to obey him,> Mary argued.

<Does tha' mean to say,> cried Martha with wide open eyes, <that he was nice to thee!>

<I think he almost liked me,> Mary answered.

<Then tha' must have bewitched him!> decided Martha, drawing a long breath.

<Do you mean Magic?> inquired Mary. <I've heard about Magic in India, but I can't make it. I just went into his room and I was so surprised to see him I stood and stared. And then he turned round and stared at me. And he thought I was a ghost or a dream and I thought perhaps he was. And it was so queer being there alone together in the middle of the night and not knowing about each other. And we began to ask each other questions. And when I asked him if I must go away he said I must not.>

<Th' world's comin' to a end!> gasped Martha.

<What is the matter with him?> asked Mary.

<Nobody knows for sure and certain,> said Martha. <Mr~Craven went off his head like when he was born. Th' doctors thought he'd have to be put in a 'sylum. It was because Mrs~Craven died like I told you. He wouldn't set eyes on th' baby. He just raved and said it'd be another hunchback like him and it'd better die.>

<Is Colin a hunchback?> Mary asked. <He didn't look like one.>

<He isn't yet,> said Martha. <But he began all wrong. Mother said that there was enough trouble and raging in th' house to set any child wrong. They was afraid his back was weak an' they've always been takin' care of it—keepin' him lyin' down and not lettin' him walk. Once they made him wear a brace but he fretted so he was downright ill. Then a big doctor came to see him an' made them take it off. He talked to th' other doctor quite rough—in a polite way. He said there'd been too much medicine and too much lettin' him have his own way.>

<I think he's a very spoiled boy,> said Mary.

<He's th' worst young nowt as ever was!> said Martha. <I won't say as he hasn't been ill a good bit. He's had coughs an' colds that's nearly killed him two or three times. Once he had rheumatic fever an' once he had typhoid. Eh! Mrs~Medlock did get a fright then. He'd been out of his head an' she was talkin' to th' nurse, thinkin' he didn't know nothin', an' she said, 'He'll die this time sure enough, an' best thing for him an' for everybody.' An' she looked at him an' there he was with his big eyes open, starin' at her as sensible as she was herself. She didn't know what'd happen but he just stared at her an' says, 'You give me some water an' stop talkin'.'>

<Do you think he will die?> asked Mary.

<Mother says there's no reason why any child should live that gets no fresh air an' doesn't do nothin' but lie on his back an' read picture-books an' take medicine. He's weak and hates th' trouble o' bein' taken out o' doors, an' he gets cold so easy he says it makes him ill.>

Mary sat and looked at the fire.

<I wonder,> she said slowly, <if it would not do him good to go out into a garden and watch things growing. It did me good.>

<One of th' worst fits he ever had,> said Martha, <was one time they took him out where the roses is by the fountain. He'd been readin' in a paper about people gettin' somethin' he called 'rose cold' an' he began to sneeze an' said he'd got it an' then a new gardener as didn't know th' rules passed by an' looked at him curious. He threw himself into a passion an' he said he'd looked at him because he was going to be a hunchback. He cried himself into a fever an' was ill all night.>

<If he ever gets angry at me, I'll never go and see him again,> said Mary.

<He'll have thee if he wants thee,> said Martha. <Tha' may as well know that at th' start.>

Very soon afterward a bell rang and she rolled up her knitting.

<I dare say th' nurse wants me to stay with him a bit,> she said. <I hope he's in a good temper.>

She was out of the room about ten minutes and then she came back with a puzzled expression.

<Well, tha' has bewitched him,> she said. <He's up on his sofa with his picture-books. He's told the nurse to stay away until six o'clock. I'm to wait in the next room. Th' minute she was gone he called me to him an' says, 'I want Mary Lennox to come and talk to me, and remember you're not to tell any one.' You'd better go as quick as you can.>

Mary was quite willing to go quickly. She did not want to see Colin as much as she wanted to see Dickon, but she wanted to see him very much.

There was a bright fire on the hearth when she entered his room, and in the daylight she saw it was a very beautiful room indeed. There were rich colours in the rugs and hangings and pictures and books on the walls which made it look glowing and comfortable even in spite of the gray sky and falling rain. Colin looked rather like a picture himself. He was wrapped in a velvet dressing-gown and sat against a big brocaded cushion. He had a red spot on each cheek.

<Come in,> he said. <I've been thinking about you all morning.>

<I've been thinking about you, too,> answered Mary. <You don't know how frightened Martha is. She says Mrs~Medlock will think she told me about you and then she will be sent away.>

He frowned.

<Go and tell her to come here,> he said. <She is in the next room.>

Mary went and brought her back. Poor Martha was shaking in her shoes. Colin was still frowning.

<Have you to do what I please or have you not?> he demanded.

<I have to do what you please, sir,> Martha faltered, turning quite red.

<Has Medlock to do what I please?>

<Everybody has, sir,> said Martha.

<Well, then, if I order you to bring Miss Mary to me, how can Medlock send you away if she finds it out?>

<Please don't let her, sir,> pleaded Martha.

<I'll send her away if she dares to say a word about such a thing,> said Master Craven grandly. <She wouldn't like that, I can tell you.>

<Thank you, sir,> bobbing a curtsy, <I want to do my duty, sir.>

<What I want is your duty,> said Colin more grandly still. <I'll take care of you. Now go away.>

When the door closed behind Martha, Colin found Mistress Mary gazing at him as if he had set her wondering.

<Why do you look at me like that?> he asked her. <What are you thinking about?>

<I am thinking about two things.>

<What are they? Sit down and tell me.>

<This is the first one,> said Mary, seating herself on the big stool. <Once in India I saw a boy who was a Rajah. He had rubies and emeralds and diamonds stuck all over him. He spoke to his people just as you spoke to Martha. Everybody had to do everything he told them—in a minute. I think they would have been killed if they hadn't.>

<I shall make you tell me about Rajahs presently,> he said, <but first tell me what the second thing was.>

<I was thinking,> said Mary, <how different you are from Dickon.>

<Who is Dickon?> he said. <What a queer name!>

She might as well tell him, she thought. She could talk about Dickon without mentioning the secret garden. She had liked to hear Martha talk about him. Besides, she longed to talk about him. It would seem to bring him nearer.

<He is Martha's brother. He is twelve years old,> she explained. <He is not like any one else in the world. He can charm foxes and squirrels and birds just as the natives in India charm snakes. He plays a very soft tune on a pipe and they come and listen.>

There were some big books on a table at his side and he dragged one suddenly toward him.

<There is a picture of a snake-charmer in this,> he exclaimed. <Come and look at it.>

The book was a beautiful one with superb coloured illustrations and he turned to one of them.

<Can he do that?> he asked eagerly.

<He played on his pipe and they listened,> Mary explained. <But he doesn't call it Magic. He says it's because he lives on the moor so much and he knows their ways. He says he feels sometimes as if he was a bird or a rabbit himself, he likes them so. I think he asked the robin questions. It seemed as if they talked to each other in soft chirps.>

Colin lay back on his cushion and his eyes grew larger and larger and the spots on his cheeks burned.

<Tell me some more about him,> he said.

<He knows all about eggs and nests,> Mary went on. < And he knows where foxes and badgers and otters live. He keeps them secret so that other boys won't find their holes and frighten them. He knows about everything that grows or lives on the moor.>

<Does he like the moor?> said Colin. <How can he when it's such a great, bare, dreary place?>

<It's the most beautiful place,> protested Mary. <Thousands of lovely things grow on it and there are thousands of little creatures all busy building nests and making holes and burrows and chippering or singing or squeaking to each other. They are so busy and having such fun under the earth or in the trees or heather. It's their world.>

<How do you know all that?> said Colin, turning on his elbow to look at her.

<I have never been there once, really,> said Mary suddenly remembering. <I only drove over it in the dark. I thought it was hideous. Martha told me about it first and then Dickon. When Dickon talks about it you feel as if you saw things and heard them and as if you were standing in the heather with the sun shining and the gorse smelling like honey—and all full of bees and butterflies.>

<You never see anything if you are ill,> said Colin restlessly. He looked like a person listening to a new sound in the distance and wondering what it was.

<You can't if you stay in a room,> said Mary.

<I couldn't go on the moor,> he said in a resentful tone.

Mary was silent for a minute and then she said something bold.

<You might—sometime.>

He moved as if he were startled.

<Go on the moor! How could I\@? I am going to die.>

<How do you know?> said Mary unsympathetically. She didn't like the way he had of talking about dying. She did not feel very sympathetic. She felt rather as if he almost boasted about it.

<Oh, I've heard it ever since I remember,> he answered crossly. <They are always whispering about it and thinking I don't notice. They wish I would, too.>

Mistress Mary felt quite contrary. She pinched her lips together.

<If they wished I would,> she said, <I wouldn't. Who wishes you would?>

<The servants—and of course Dr~Craven because he would get Misselthwaite and be rich instead of poor. He daren't say so, but he always looks cheerful when I am worse. When I had typhoid fever his face got quite fat. I think my father wishes it, too.>

<I don't believe he does,> said Mary quite obstinately.

That made Colin turn and look at her again.

<Don't you?> he said.

And then he lay back on his cushion and was still, as if he were thinking. And there was quite a long silence. Perhaps they were both of them thinking strange things children do not usually think of.

<I like the grand doctor from London, because he made them take the iron thing off,> said Mary at last. <Did he say you were going to die?>

<No.>

<What did he say?>

<He didn't whisper,> Colin answered. <Perhaps he knew I hated whispering. I heard him say one thing quite aloud. He said, 'The lad might live if he would make up his mind to it. Put him in the humor.' It sounded as if he was in a temper.>

<I'll tell you who would put you in the humor, perhaps,> said Mary reflecting. She felt as if she would like this thing to be settled one way or the other. <I believe Dickon would. He's always talking about live things. He never talks about dead things or things that are ill. He's always looking up in the sky to watch birds flying—or looking down at the earth to see something growing. He has such round blue eyes and they are so wide open with looking about. And he laughs such a big laugh with his wide mouth—and his cheeks are as red—as red as cherries.>

She pulled her stool nearer to the sofa and her expression quite changed at the remembrance of the wide curving mouth and wide open eyes.

<See here,> she said. <Don't let us talk about dying; I don't like it. Let us talk about living. Let us talk and talk about Dickon. And then we will look at your pictures.>

It was the best thing she could have said. To talk about Dickon meant to talk about the moor and about the cottage and the fourteen people who lived in it on sixteen shillings a week—and the children who got fat on the moor grass like the wild ponies. And about Dickon's mother—and the skipping-rope—and the moor with the sun on it—and about pale green points sticking up out of the black sod. And it was all so alive that Mary talked more than she had ever talked before—and Colin both talked and listened as he had never done either before. And they both began to laugh over nothings as children will when they are happy together. And they laughed so that in the end they were making as much noise as if they had been two ordinary healthy natural ten-year-old creatures—instead of a hard, little, unloving girl and a sickly boy who believed that he was going to die.

They enjoyed themselves so much that they forgot the pictures and they forgot about the time. They had been laughing quite loudly over Ben Weatherstaff and his robin and Colin was actually sitting up as if he had forgotten about his weak back when he suddenly remembered something.

<Do you know there is one thing we have never once thought of,> he said. <We are cousins.>

It seemed so queer that they had talked so much and never remembered this simple thing that they laughed more than ever, because they had got into the humor to laugh at anything. And in the midst of the fun the door opened and in walked Dr~Craven and Mrs~Medlock.

Dr~Craven started in actual alarm and Mrs~Medlock almost fell back because he had accidentally bumped against her.

<Good Lord!> exclaimed poor Mrs~Medlock, with her eyes almost starting out of her head. <Good Lord!>

<What is this?> said Dr~Craven, coming forward. <What does it mean?>

Then Mary was reminded of the boy Rajah again. Colin answered as if neither the doctor's alarm nor Mrs~Medlock's terror were of the slightest consequence. He was as little disturbed or frightened as if an elderly cat and dog had walked into the room.

<This is my cousin, Mary Lennox,> he said. <I asked her to come and talk to me. I like her. She must come and talk to me whenever I send for her.>

Dr~Craven turned reproachfully to Mrs~Medlock.

<Oh, sir,> she panted. <I don't know how it's happened. There's not a servant on the place that'd dare to talk—they all have their orders.>

<Nobody told her anything,> said Colin, <she heard me crying and found me herself. I am glad she came. Don't be silly, Medlock.>

Mary saw that Dr~Craven did not look pleased, but it was quite plain that he dare not oppose his patient. He sat down by Colin and felt his pulse.

<I am afraid there has been too much excitement. Excitement is not good for you, my boy,> he said.

<I should be excited if she kept away,> answered Colin, his eyes beginning to look dangerously sparkling. <I am better. She makes me better. The nurse must bring up her tea with mine. We will have tea together.>

Mrs~Medlock and Dr~Craven looked at each other in a troubled way, but there was evidently nothing to be done.

<He does look rather better, sir,> ventured Mrs~Medlock. <But>—thinking the matter over—<he looked better this morning before she came into the room.>

<She came into the room last night. She stayed with me a long time. She sang a Hindustani song to me and it made me go to sleep,> said Colin. <I was better when I wakened up. I wanted my breakfast. I want my tea now. Tell nurse, Medlock.>

Dr~Craven did not stay very long. He talked to the nurse for a few minutes when she came into the room and said a few words of warning to Colin. He must not talk too much; he must not forget that he was ill; he must not forget that he was very easily tired. Mary thought that there seemed to be a number of uncomfortable things he was not to forget.

Colin looked fretful and kept his strange black-lashed eyes fixed on Dr~Craven's face.

<I want to forget it,> he said at last. <She makes me forget it. That is why I want her.>

Dr~Craven did not look happy when he left the room. He gave a puzzled glance at the little girl sitting on the large stool. She had become a stiff, silent child again as soon as he entered and he could not see what the attraction was. The boy actually did look brighter, however—and he sighed rather heavily as he went down the corridor.

<They are always wanting me to eat things when I don't want to,> said Colin, as the nurse brought in the tea and put it on the table by the sofa. <Now, if you'll eat I will. Those muffins look so nice and hot. Tell me about Rajahs.>